% SPDX-License-Identifier: Apache-2.0
% covers: EthTx
\datasheet{EthTx}{The transmit half of the Ethernet MAC: whole frames from a channel onto GMII.}

\dsfacts{\code{lib/parts/src/eth.rs}}{\code{txhdl\_parts::eth::EthTx}}%
{\code{//docs:eth}}

\subsection*{Function}

\code{EthTx} takes a frame's bytes from the channel \code{tx}, each an
\code{EthByte}: the byte and whether it is the frame's last. A frame
there is the destination address, the source address, the type or
length and the payload. The unit stores the whole frame, then sends it
on GMII: \code{txd}, a byte per cycle, and \code{tx\_en}. It adds the
preamble, the start of frame delimiter, padding up to sixty bytes and
the four bytes of the CRC-32 frame check sequence. The unit runs on one
clock, so on a board it can run on the transmit clock.

\subsection*{Parameters}

\code{EthTx} has no type parameters. Its frame memory holds
\code{FRAME\_MAX} bytes, 2048, a constant of the module.

\dstables{EthTx}

\subsection*{Behaviour}

The unit is two processes. The store takes the frame in; the sender is
the wire's protocol written as the sequence it is: wait for a whole
frame, the preamble and the delimiter, the frame's bytes, zeros up to
sixty, the check sequence and the gap, a byte a cycle, each a turn of
a loop. The lowering numbers the sender's waits into a state register
the unit does not declare, \code{at\_wait}, keeps its loops' counts in
\code{for0} to \code{for3}, and holds the two outputs between states in
\code{txd\_held} and \code{tx\_en\_held}; all are in the tables above.

\begin{itemize}
\item While no frame is whole, the store takes a byte from \code{tx}
whenever one is offered, and stores it in \code{frame} at \code{fill}.
The byte marked last sets \code{whole}.
\item A whole frame starts the sender: seven bytes of \code{0x55}, then
\code{0xd5}, with the CRC register set to \code{CRC\_INIT}, every bit
set. The frame follows, the register running over it.
\item A frame of fewer than sixty bytes is followed by zeros up to
sixty, with the register still running.
\item The check sequence is the register complemented, least
significant byte first.
\item \code{tx\_en} is high from the preamble to the end of the check
sequence. Twelve bytes of gap follow with \code{tx\_en} low. Then
\code{sent} is up for one cycle, which adds one to \code{frames} and
lets the store forget the frame and take the next.
\item The store takes no byte from \code{tx} while a frame is whole,
from its last byte stored to the end of its gap, so it holds the client
off. The client may offer bytes as slowly as it likes.
\item \code{crc\_byte} steps the register by a whole byte: the byte
added to it, then eight shifts with the polynomial \code{0xedb88320}
going in whenever the bit that leaves is set.
\end{itemize}

\subsection*{Verification}

\begin{itemize}
\item The tests below are in \code{lib/parts/src/eth.rs}, and run in
\code{//lib/parts:parts\_test}.
\item \code{the\_byte\_wide\_crc\_is\_the\_bitwise\_one} checks
\code{crc\_byte} against the register stepped a bit at a time, on every
byte from four registers.
\item \code{the\_crc\_of\_a\_known\_string} checks \code{crc32} against
the catalogue value and the residue.
\item \code{a\_frame\_goes\_out\_as\_the\_model\_says\_and\_comes\_back}
loops the transmitter into \code{EthRx}, and holds the bytes on the
wire to \code{wire\_bytes}, for frames of 20 and 100 bytes.
\item \code{frames\_at\_the\_minimum\_length} does the same for frames
of 59, 60 and 61 bytes.
\item \code{ex\_eth} runs the unit behind \code{EthLite} with the wire
looped back, and asserts that it sent three frames. The build simulates
the netlist of \code{EthTx} against that run under nvc and Verilator:
\code{//docs:eth\_sim\_eth\_tx\_eth\_tx\_tb\_test} and
\code{//docs:eth\_vsim\_eth\_tx\_test}.
\end{itemize}

\subsection*{Used by}

The example \code{ex\_eth}. The echo design for the AX7A200B:
\code{//eth:netlist} writes the Verilog of \code{eth\_tx}, and
\code{eth/board/eth\_echo.v} instantiates it on the 125~MHz transmit
clock. \code{//eth:echo\_synth} and \code{//eth:echo\_pnr} put that
design through Vivado to a bitstream. \code{//docs:eth} states that it
has not yet run on the board.

\subsection*{Limits}

The unit holds one frame. A frame longer than \code{FRAME\_MAX} bytes
keeps overwriting its last byte. Only gigabit works: at 100 or
10~Mbit a PHY sends a nibble per cycle, and the unit does not take a
byte over two cycles.
